\section{Gauge constraints, observable algebras, and physical Hilbert spaces}
\label{sec:gauge-aqft}
% BEGIN CURATED SUBJECT INDEX
\index{Hilbert space}
\index{Gauss law}
\index{BRST cohomology}
\index{observable algebra}
\index{algebraic quantum field theory}
% END CURATED SUBJECT INDEX

Long-range QED exposes a question that is already present in every gauge
theory: which Hilbert space is physical?  A gauge-fixed field algebra, a
BRST complex, a gauge-invariant observable algebra, and the Hilbert space used
by a lattice simulator are related, but they are not the same object.  Direct
integrals clarify the relation only if they are applied at the correct stage.
They decompose a representation that has already been selected; they do not
select the physical representations by themselves.

This section gives an observable-first construction and compares it with the
BRST route.  It then separates two sources of nonfactorization: generic local
QFT structure and the additional constraint imposed by Gauss law.  The aim is
not a nonperturbative construction of Yang--Mills theory.  It is a precise
list of the data such a construction would have to provide.

\subsection{Three layers that must not be collapsed}
% BEGIN CURATED SUBJECT INDEX
\index{Hilbert space}
\index{GNS construction}
\index{observable algebra}
\index{algebraic quantum field theory}
% END CURATED SUBJECT INDEX

It is useful to keep three mathematical layers visible:
\begin{enumerate}[label=(\roman*)]
  \item a kinematical field space, possibly gauge fixed and possibly carrying
  an indefinite inner product;
  \item an abstract net or algebra $\vsub{\Alg}{obs}$ of gauge-invariant
  observables; and
  \item one or more positive Hilbert-space representations
  $(\pi,\Hil,\Omega)$ selected by a state, an energy condition, and a charge
  or asymptotic criterion.
\end{enumerate}
Gauge-dependent fields can be invaluable coordinates for calculation, while
the observable algebra records what can be measured without choosing those
coordinates.  A Hilbert representation adds the state-dependent meaning of
expectation values, spectrum, and weak closure.  In systems with infinitely
many degrees of freedom there need not be one representation containing all
physically relevant sectors.

For a relativistic theory the observable-first datum is a net
\begin{equation}
  \mathcal O\longmapsto\Alg(\mathcal O),
  \qquad
  \mathcal O_1\subset\mathcal O_2
  \Longrightarrow
  \Alg(\mathcal O_1)\subset\Alg(\mathcal O_2),
  \label{eq:aqft-isotony}
\end{equation}
with locality
\begin{equation}
  [A_1,A_2]=0
  \quad\text{if }\mathcal O_1\text{ and }\mathcal O_2
  \text{ are spacelike separated}.
  \label{eq:aqft-locality}
\end{equation}
Covariance, a spectrum condition, and an appropriate vacuum or thermal state
complete the basic data.  The quasilocal $C^*$ algebra is obtained by norm
completion of the union of local algebras.  A state $\omega$ then constructs
the Hilbert space through GNS:
\begin{equation}
  \omega(A)
  =\dualpair{\Omega_\omega}{
  \pi_\omega(A)\Omega_\omega},
  \qquad
  \Hil_\omega
  =\overline{\pi_\omega(\vsub{\Alg}{obs})\Omega_\omega}.
  \label{eq:gauge-observable-gns}
\end{equation}
Thus positivity belongs to the state on observables, not to a particular
gauge potential.

\subsection{The BRST cohomological route}
% BEGIN CURATED SUBJECT INDEX
\index{operator!closed}
\index{Hilbert space}
\index{topology}
\index{BRST cohomology}
\index{Gribov problem}
% END CURATED SUBJECT INDEX

Covariant gauge fixing enlarges the field space by longitudinal modes,
ghosts, and auxiliary fields.  Let $\mathcal K$ denote the resulting Krein
space and $Q$ a nilpotent BRST charge,
\begin{equation}
  Q^2=0.
  \label{eq:brst-nilpotent}
\end{equation}
The algebraic physical space is the cohomology
\begin{equation}
  \vsupb{\Hil}{alg}{BRST}
  =\ker Q/\operatorname{im}Q.
  \label{eq:brst-cohomology}
\end{equation}
To obtain a Hilbert space one still has to prove that the induced sesquilinear
form is positive semidefinite, quotient its null vectors, and complete the
result.  Observables are represented by BRST-closed operators modulo exact
ones.  In the Kugo--Ojima framework this implements the quartet mechanism and
connects color confinement to the global BRST charge under explicit
assumptions \cite{KugoOjima:1979}.

Equation~\eqref{eq:brst-cohomology} is elegant because local covariant fields
and positivity are reconciled after taking cohomology.  It is not, however, a
definition that can simply be assumed in an arbitrary nonperturbative theory.
One needs a globally defined nilpotent charge, a controlled domain, positivity
of the quotient, and a state space on which the construction is complete.
Gribov copies obstruct a global gauge slice, and their topology warns that a
perturbative BRST complex need not automatically extend to the full
configuration space \cite{Gribov:1978,Singer:1978}.  Abandoning these global
questions may be a deliberate approximation, but it must then be identified
as part of the truncation contract.

\subsection{The observable-first definition}
% BEGIN CURATED SUBJECT INDEX
\index{Hilbert space}
\index{measure theory}
\index{direct integral}
\index{factor}
\index{GNS construction}
\index{BRST cohomology}
\index{observable algebra}
% END CURATED SUBJECT INDEX

If BRST cohomology is not taken as the global starting point, a nonperturbative
physical Hilbert space should be defined in the following order:
\begin{enumerate}[label=\arabic*.]
  \item construct a positive gauge-invariant observable algebra or net,
  including its dynamics;
  \item select a class of physical states by positivity, the Gauss constraint,
  energy and regularity conditions, and the desired vacuum, thermal, charge,
  or asymptotic behavior;
  \item apply GNS to each selected state or folium;
  \item compare the resulting representations by unitary equivalence,
  quasi-equivalence, or disjointness; and
  \item only then form direct sums or direct integrals when a specified mixed
  representation and measure call for them.
\end{enumerate}

In symbols, a selected family $\{\omega_\xi\}_{\xi\in X}$ can produce
\begin{equation}
  \pi_\mu
  =\int_X^\oplus\pi_\xi\dd\mu(\xi),
  \qquad
  \Hil_\mu
  =\int_X^\oplus\Hil_\xi\dd\mu(\xi),
  \label{eq:physical-sector-direct-integral}
\end{equation}
provided the fields are measurable.  The measure $\mu$ is part of the state
or ensemble; it is not determined by the abstract algebra alone.  A pure
sector may require only one factor representation.  A finite compact lattice
problem more often gives a direct \textit{sum}.  A genuine direct integral is
natural for continuous charge or flux data, infinite volume, or a central
decomposition with a nonatomic center.

This answers the apparent dilemma.  Insisting on a direct integral ``on the
observable side'' is legitimate as a posterior decomposition of a chosen
representation.  It cannot replace the state-selection problem.  Without the
selected state class and measure, the phrase ``the physical Hilbert space''
is underdetermined.

\subsection{Central decomposition is not sector classification}
% BEGIN CURATED SUBJECT INDEX
\index{von Neumann algebra@\vN\ algebra}
\index{central decomposition}
\index{superselection sector}
% END CURATED SUBJECT INDEX

Let
$\vNAlg_\omega=\pi_\omega(\vsub{\Alg}{obs})''$.  Reduction theory gives,
under the hypotheses of Section~\ref{subsec:central-decomposition},
\begin{equation}
  \vNAlg_\omega
  =\int_X^\oplus\vNAlg_{\omega,\xi}\dd\mu_\omega(\xi),
  \qquad
  \Center(\vNAlg_{\omega,\xi})
  =\complex\identity_\xi
  \quad\text{a.e.}
  \label{eq:gauge-central-decomposition}
\end{equation}
This decomposes one represented algebra into factors.  A superselection
theory asks a different question: which representations of the abstract net
are physically admissible, and how do charges compose?  The DHR criterion
selects charges indistinguishable from the vacuum outside bounded regions
\cite{DoplicherHaagRoberts:1971}.  Gauss-law charges are detectable at
arbitrarily large distances and do not satisfy that compact-localization
criterion.  Spacelike-cone localization and asymptotic constructions provide
weaker alternatives in suitable theories
\cite{BuchholzFredenhagen:1982,Buchholz:1986}.

Consequently, a factor fiber in
Eq.~\eqref{eq:gauge-central-decomposition} should be called a factor
component until a physical selection criterion identifies it with a charge
sector.  Conversely, two charge sectors may be disjoint representations of
the same simple quasilocal algebra even though the abstract algebra has
trivial center.  The center appears after representation or after restricting
to a region with boundary data; it need not be present in the global abstract
algebra.

\subsection{Solving Gauss law on a lattice}
% BEGIN CURATED SUBJECT INDEX
\index{Gauss law}
\index{tensor product!gauge constraint}
% END CURATED SUBJECT INDEX

Hamiltonian lattice gauge theory makes the reorganization operational
\cite{KogutSusskind:1975}.  Before imposing constraints, a finite graph with
links $e$ has a kinematical tensor product
\begin{equation}
  \vsub{\Hil}{kin}
  =\bigotimes_e\Hil_e
  \otimes\bigotimes_v\vsup{\Hil_v}{matter}.
  \label{eq:lattice-kinematical-tensor-product}
\end{equation}
At every vertex, the Gauss operator $G_v$ generates a gauge transformation.
The physical subspace is
\begin{equation}
  \vsub{\Hil}{phys}
  =\{\psi\in\vsub{\Hil}{kin}:
    \text{$G_v\psi=\psi$ for all $v$}\}
  \label{eq:lattice-gauss-subspace}
\end{equation}
for a finite gauge group, with the corresponding zero-generator condition
for a Lie group.  The orthogonal projector is the group average
\begin{equation}
  P_G=\prod_v\int_G\rmd g_v\,U_v(g_v).
  \label{eq:lattice-gauss-projector}
\end{equation}
Since a vertex transformation acts on all incident links,
$P_G$ does not factor into independent projectors for arbitrary spatial
subregions.

One can solve the constraints explicitly by choosing a spanning tree and
expressing tree-link variables through matter charges and the remaining loop
variables.  This can reduce the number of degrees of freedom and is valuable
for quantum simulation \cite{PardoEtAl:2023}.  The price is that an operator
local in the original link variables may become a string along the chosen
tree.  The exact shape of that string is encoding dependent.  The need to
carry boundary flux and the failure of a naive regional tensor product are
not.

\subsection{Laboratory: the smallest \texorpdfstring{$\mathbb Z_2$}{Z2}
boundary-flux sector}
% BEGIN CURATED SUBJECT INDEX
\index{Hilbert space}
\index{direct integral}
\index{central decomposition}
\index{Gauss law}
\index{observable algebra}
\index{tensor product!gauge constraint}
% END CURATED SUBJECT INDEX

Consider a connected region whose interior is a single link joining two
vertices.  One external link crosses the boundary at each vertex.  In the
electric basis write the $\mathbb Z_2$ labels as
$q_L$, $q$, $q_R\in\{0,1\}$.  With no interior matter charge, Gauss law gives
\begin{equation}
  q_L+q=0,
  \qquad
  q+q_R=0
  \pmod 2.
  \label{eq:z2-two-vertex-gauss}
\end{equation}
Hence $q_L=q=q_R$.  If the exterior is not included in the regional Hilbert
space, the physical regional space has the decomposition
\begin{equation}
  \vsup{\Hil_R}{phys}
  =\Hil_{R,0}\oplus\Hil_{R,1}.
  \label{eq:z2-boundary-direct-sum}
\end{equation}
Each summand has fixed boundary electric flux.  Gauge-invariant operators
strictly inside the region cannot change that label.  The boundary-flux
operator therefore lies in the center of the regional observable algebra and
the projectors onto $\Hil_{R,0}$ and $\Hil_{R,1}$ are central projections.

For this finite compact example the correct structure is a direct sum, not a
direct integral.  Joining the region to its exterior requires matching the
left and right boundary representations and contracting their indices.  An
``extended Hilbert space'' can instead introduce edge variables so that a
tensor product is restored before imposing the matching condition
\cite{CasiniHuertaRosabal:2014,DonnellyFreidel:2016}.  These are different
encodings of the same gauge-invariant content, provided the boundary algebra
and its center are tracked.

\subsection{Which nonlocality is an artifact?}
% BEGIN CURATED SUBJECT INDEX
\index{von Neumann algebra@\vN\ algebra}
\index{superselection sector}
\index{dressing state}
\index{Gauss law}
\index{observable algebra}
\index{algebraic quantum field theory}
\index{tensor product!gauge constraint}
% END CURATED SUBJECT INDEX

The answer has two parts.
\begin{description}[style=nextline]
  \item[Encoding-dependent nonlocality.]
  A spanning tree, Jordan--Wigner-like ordering, gauge fixing, or choice of
  edge modes determines where strings are drawn and which qubits store a loop
  variable.  Changing the encoding can move or shorten those strings.

  \item[Structural nonlocality or nonfactorization.]
  Gauss law ties charge in a region to flux through its boundary.  Charged
  fields cannot be both gauge invariant and compactly localized without an
  accompanying dressing.  Regional observable algebras possess boundary
  superselection data, and the physical subspace generally fails to equal a
  tensor product of independently physical subregions.
\end{description}
Thus not every long Pauli string seen by a simulator is fundamental, but the
constraint that makes some boundary or dressing data unavoidable is
fundamental.

There is also a logically separate AQFT effect.  Local \vN\ algebras
in continuum relativistic QFT are typically type III and therefore are not
ordinary tensor factors with density matrices, even in a theory without gauge
charge.  When nested regions $\mathcal O_1\Subset\mathcal O_2$ satisfy the
\term{split property}, there exists a type I factor $\mathcal N$ such that
\begin{equation}
  \Alg(\mathcal O_1)''
  \subset\mathcal N
  \subset\Alg(\mathcal O_2)''.
  \label{eq:split-inclusion}
\end{equation}
The collar between the regions permits an approximate subsystem structure
\cite{FewsterSplit:2015}.  Type III nonfactorization, failure of the split
property, and a Gauss-law boundary center should therefore be diagnosed
separately; calling all three ``nonlocal encoding'' loses the distinction.

\subsection{Connection coefficients sector by sector}
% BEGIN CURATED SUBJECT INDEX
\index{Hilbert space}
\index{topology}
\index{measure theory}
\index{direct integral}
\index{central decomposition}
\index{GNS construction}
\index{connection coefficient}
\index{boundary condition!incoming and outgoing}
\index{resonance!residue}
\index{dressing state}
% END CURATED SUBJECT INDEX

The observable-first picture meets the analytic strategy of this book when a
Hamiltonian or radial problem decomposes over a measurable sector variable:
\begin{equation}
  \Hil_\mu=\int_X^\oplus\Hil_\xi\dd\mu(\xi),
  \qquad
  H=\int_X^\oplus H_\xi\dd\mu(\xi).
  \label{eq:sectorwise-hamiltonian}
\end{equation}
If each fiber has a connection problem, exact WKB produces an incoming
coefficient $\Cincoming(E,\xi)$.  A fiber resonance satisfies
\begin{equation}
  \Cincoming(\resonantE(\xi),\xi)=0.
  \label{eq:sectorwise-resonance-condition}
\end{equation}
But performing this calculation independently in every fiber is not yet an
operator theorem.  One must prove measurability in $\xi$, choose analytic
continuation domains uniform on the relevant set of sectors, control residues
near thresholds, and justify interchanging continuation with the direct
integral.  If pole curves collide or escape through a sector-dependent cut,
no single contour need isolate them uniformly.

This is a concrete open interface between exact WKB and operator algebras.
The connection coefficient supplies analytic information that the measurable
decomposition alone lacks; the direct integral supplies the topology and
measure needed to assemble fiber calculations into a physical state space.
Neither construction subsumes the other.

\begin{checkpointbox}
A nonperturbative gauge theory is not assigned a physical Hilbert space by
writing a direct integral.  It is assigned an observable algebra, dynamics,
and a physically selected state class; GNS constructs the representations.
Direct sums, direct integrals, and central decompositions then organize those
representations.  Solving Gauss law can make a chosen encoding nonlocal, while
boundary flux, charged dressing, and nonfactorization are structural data that
any faithful encoding must preserve.
\end{checkpointbox}
